% ============================================================
%  PATENT APPLICATION
%  System and Method for Offline Semantic Pre-Computation
%  and Knowledge Graph Construction from Timestamped Text Corpora
% ============================================================
\documentclass[12pt,letterpaper]{article}

% ---- Packages ----------------------------------------------
\usepackage[top=1.25in, bottom=1.25in, left=1.25in, right=1.25in]{geometry}
\usepackage{fontenc}
\usepackage{inputenc}
\usepackage{lmodern}
\usepackage{microtype}
\usepackage{setspace}
\usepackage{titlesec}
\usepackage{enumitem}
\usepackage{booktabs}
\usepackage{longtable}
\usepackage{array}
\usepackage{parskip}
\usepackage{fancyhdr}
\usepackage{listings}
\usepackage{xcolor}
\usepackage{mdframed}
\usepackage{tocloft}
\usepackage{hyperref}

% ---- Colors ------------------------------------------------
\definecolor{codebg}{gray}{0.96}
\definecolor{codeframe}{gray}{0.75}
\definecolor{keywordcolor}{rgb}{0.0, 0.3, 0.7}
\definecolor{commentcolor}{rgb}{0.3, 0.5, 0.3}
\definecolor{stringcolor}{rgb}{0.6, 0.1, 0.1}

% ---- Listings setup ----------------------------------------
\lstset{
  basicstyle=\ttfamily\small,
  backgroundcolor=\color{codebg},
  frame=single,
  framerule=0.6pt,
  rulecolor=\color{codeframe},
  breaklines=true,
  breakatwhitespace=false,
  keepspaces=true,
  showspaces=false,
  showstringspaces=false,
  tabsize=2,
  captionpos=b,
  numbers=none,
  xleftmargin=6pt,
  xrightmargin=6pt,
}

\lstdefinelanguage{YAML}{
  keywords={true,false,null,~},
  sensitive=false,
  comment=[l]{\#},
  morestring=[b]',
  morestring=[b]",
  keywordstyle=\color{keywordcolor}\bfseries,
  commentstyle=\color{commentcolor}\itshape,
  stringstyle=\color{stringcolor},
}

\lstdefinelanguage{JSON}{
  morestring=[b]",
  stringstyle=\color{stringcolor},
  literate=
    *{:}{{{\color{black}{:}}}}{1}
    {,}{{{\color{black}{,}}}}{1}
    {\{}{{{\color{black}{\{}}}}{1}
    {\}}{{{\color{black}{\}}}}}{1}
    {[}{{{\color{black}{[}}}}{1}
    {]}{{{\color{black}{]}}}}{1},
}

% ---- Hyperref setup ----------------------------------------
\hypersetup{
  colorlinks=false,
  pdfborder={0 0 0},
  pdftitle={Patent Application -- Offline Semantic Pre-Computation and Knowledge Graph Construction},
  pdfauthor={Eric G. Suchanek, PhD},
}

% ---- Section title formatting ------------------------------
\titleformat{\section}
  {\normalfont\large\bfseries\uppercase}
  {}
  {0pt}
  {}
  [\vspace{2pt}\hrule\vspace{4pt}]

\titleformat{\subsection}
  {\normalfont\normalsize\bfseries}
  {}
  {0pt}
  {}

\titlespacing*{\section}{0pt}{18pt}{8pt}
\titlespacing*{\subsection}{0pt}{12pt}{4pt}

% ---- Page style --------------------------------------------
\pagestyle{fancy}
\fancyhf{}
\fancyhead[L]{\small\textit{Patent Application}}
\fancyhead[R]{\small\textit{Suchanek -- Offline Semantic Pre-Computation}}
\fancyfoot[C]{\small\thepage}
\renewcommand{\headrulewidth}{0.4pt}

% ---- Claim formatting helper --------------------------------
\newcounter{claimnum}
\newenvironment{claims}
  {\begin{list}{\textbf{Claim \arabic{claimnum}.}}
      {\usecounter{claimnum}
       \setlength{\labelwidth}{1.8cm}
       \setlength{\labelsep}{0.4cm}
       \setlength{\leftmargin}{2.2cm}
       \setlength{\itemsep}{12pt}
       \setlength{\parsep}{4pt}}}
  {\end{list}}

% ---- Indented sub-list for claim elements ------------------
\newlist{claimelems}{enumerate}{2}
\setlist[claimelems,1]{
  label=(\alph*),
  leftmargin=1.6cm,
  itemsep=4pt,
  parsep=2pt,
}
\setlist[claimelems,2]{
  label=(\roman*),
  leftmargin=1.6cm,
  itemsep=2pt,
  parsep=1pt,
}

% ---- Document spacing --------------------------------------
\setstretch{1.15}

% ============================================================
\begin{document}

% ---- Title block -------------------------------------------
\begin{center}
  {\LARGE\bfseries PATENT APPLICATION}\\[14pt]
  {\large\bfseries SYSTEM AND METHOD FOR OFFLINE SEMANTIC\\
   PRE-COMPUTATION AND KNOWLEDGE GRAPH CONSTRUCTION\\
   FROM TIMESTAMPED TEXT CORPORA}\\[20pt]
  \begin{tabular}{ll}
    \textbf{Inventor:}        & Eric G. Suchanek, PhD \\[2pt]
    \textbf{Assignee:}        & Flux-Frontiers \\[2pt]
    \textbf{Filing Date:}     & March 26, 2026 \\[2pt]
    \textbf{Application Type:} & Utility Patent \\
  \end{tabular}
\end{center}

\vspace{18pt}
\noindent\rule{\linewidth}{0.8pt}
\vspace{6pt}

% ============================================================
\section*{Cross-Reference to Related Applications}

Not applicable.

% ============================================================
\section*{Field of the Invention}

The present invention relates generally to natural language processing and knowledge
graph construction, and more particularly to a system and method for offline semantic
pre-computation of timestamped text corpora that eliminates query-time inference while
preserving full semantic and temporal search capability.

% ============================================================
\section*{Background of the Invention}

\subsection*{The Prior Art and Its Deficiencies}

Modern semantic search systems operate under two foundational assumptions that constrain
their deployment, cost, and privacy characteristics.

\textbf{First assumption: semantic understanding requires live inference.} When a user
queries a document corpus, conventional systems invoke a large language model (LLM) at
query time --- either via cloud API or local GPU inference --- to process the query and
generate semantically relevant results. This architecture imposes per-query costs that
scale linearly with query volume, requires persistent network connectivity for
cloud-based systems, and demands specialized hardware (GPU clusters) for local
deployment. Consumer-grade hardware cannot serve these workloads at interactive
latencies for corpora exceeding trivial sizes.

\textbf{Second assumption: temporal grounding requires LLM extraction.} Dates, time
references, and temporal relationships embedded in natural language prose are
conventionally extracted by language models at either ingest or query time. For
locally-deployed small-parameter models (e.g., 4 billion parameters), this extraction
fails systematically --- such models hallucinate dates, confuse temporal formats,
produce malformed output, and cannot reliably distinguish relative from absolute time
references. This failure mode blocks local, offline execution of any temporally-aware
memory or search system.

Prior art systems addressing semantic search over document corpora include
Retrieval-Augmented Generation (RAG) pipelines, which embed documents into vector
stores but require LLM inference at query time for answer generation; traditional
knowledge graph systems, which provide structural querying but lack integrated semantic
vector search; and embedding-based search engines, which provide semantic similarity
but lack temporal provenance, topic classification, and structural graph relationships.

No existing system provides, in a single integrated pipeline executable on consumer
hardware without network connectivity: (a) multi-phase NLP enrichment with configurable
topic and context classification; (b) parallel multi-process vector embedding with no
external service dependency; (c) a hybrid knowledge graph combining structural indices
with semantic vector indices; (d) temporal provenance on every node written directly
from source metadata without LLM extraction; and (e) sub-three-minute full-corpus
ingestion at the scale of thousands of entries.

% ============================================================
\section*{Summary of the Invention}

The present invention provides a system and method for constructing a semantically
enriched, temporally grounded knowledge graph from timestamped text corpora through
offline pre-computation. The core innovation is a shift from query-time inference to
ingest-time pre-computation: all semantic understanding --- topic labeling, context
classification, vector embedding, and graph edge construction --- is computed once
during corpus ingestion. The resulting knowledge graph is a frozen, versioned,
queryable artifact that serves all subsequent queries through vector lookup and graph
traversal alone, with no model loading or inference at query time.

The system comprises six sequential processing stages orchestrated by a unified
pipeline controller:

\begin{enumerate}[leftmargin=2cm, itemsep=6pt]

\item A multi-phase NLP transformer that parses timestamped source entries, applies
configurable sentence-group chunking, performs multi-label topic classification and
single-label context classification using domain-configurable YAML vocabularies, and
supports temporally-diverse sampling;

\item A corpus emitter that writes each semantically enriched chunk as a structured
document with metadata frontmatter, embedding temporal provenance directly from parsed
source timestamps without any language model extraction;

\item A structural graph builder that indexes all chunks into a SQLite database with
BM25 lexical search capability and typed graph edges;

\item A semantic vector index builder that generates dense vector embeddings and stores
them in a columnar vector database for approximate nearest-neighbor search;

\item A multi-process parallel embedder that shards the corpus across CPU worker
processes, each loading an independent model instance with no shared state and no GIL
contention; and

\item A JSON embedding cache that stores aligned arrays of embeddings, texts, and
timestamps for downstream analysis.

\end{enumerate}

The invention inverts the conventional cost model: ingestion cost is medium and paid
once; query cost is negligible and does not increase with query volume; privacy is
preserved because no data leaves the local machine; and offline operation is fully
supported.

% ============================================================
\section*{Detailed Description of the Preferred Embodiments}

\subsection*{1.\ \ System Overview}

Referring now to the invention in detail, the system operates on timestamped text
corpora in a pipe-delimited enriched format:

\begin{lstlisting}
TIMESTAMP | TYPE | CATEGORY | CONTENT
\end{lstlisting}

Each entry carries a temporal anchor (an ISO-format timestamp), a semantic type label,
and a context category. This format is produced by a preprocessing layer that enriches
raw text with natural language processing --- named-entity recognition, sentence
segmentation, and domain-specific topic classification via configurable vocabularies.

The system has been reduced to practice and validated on a corpus of 6,450 diary
entries spanning 9.6 years (1660--1669), comprising 23,235 enriched lines totaling
3.3 megabytes. Full pipeline execution completes in approximately three minutes on
consumer-grade Apple Silicon hardware with no GPU and no network connectivity.

% -----------------------------------------------------------
\subsection*{2.\ \ Stage 1: Multi-Phase NLP Transformation}

The first processing stage applies five sequential NLP phases to the enriched source
text:

\textbf{Phase 1 --- Parsing.} The system reads pipe-delimited lines and validates ISO
timestamps. Malformed entries are logged and excluded. Each valid entry produces a
structured record with timestamp, type, category, and content fields.

\textbf{Phase 2 --- Sentence-Group Chunking.} The content of each entry is segmented
into chunks using a configurable chunking strategy. The preferred embodiment uses
sentence-group chunking with a default of four sentences per chunk and a maximum of 512
characters. Alternative strategies include hybrid chunking (combining sentence
boundaries with character limits) and semantic chunking (grouping sentences by topical
coherence). The chunking parameters are configurable to accommodate different corpus
characteristics.

\textbf{Phase 3 --- Multi-Label Topic Classification.} Each chunk is classified against
a domain-configurable YAML vocabulary that maps topic labels to keyword lists.
Classification uses TF-IDF weighted keyword matching. Each chunk may carry multiple
topic labels (e.g., a single chunk may be labeled with both ``domestic'' and ``court''
topics). The vocabulary is fully user-configurable, enabling domain adaptation without
code changes.

\textbf{Phase 4 --- Single-Label Context Classification.} Each chunk receives exactly
one context label from a second configurable YAML vocabulary. Context categories
represent the situational frame of the content (e.g., ``Office,'' ``Home,'' ``Social,''
``Reflection,'' ``Health''). Assignment uses keyword matching with priority ordering
defined in the vocabulary.

\textbf{Phase 5 --- Temporally-Diverse Sampling.} When a batch size is specified (for
processing a subset of a large corpus), the system applies temporally-diverse sampling
using uniform stride across the time axis rather than random sampling. This ensures
representative coverage across the full temporal span of the corpus, preventing
front-loading bias that would result from naive sequential or random sampling.

\textbf{Parallelism.} The NLP transformation stage accepts a configurable
\texttt{workers} parameter that distributes feature extraction across CPU cores via
process-level parallelism. Each worker processes an independent shard of the corpus
with no shared state, enabling near-linear scaling with available CPU cores.

% -----------------------------------------------------------
\subsection*{3.\ \ Stage 2: Corpus Emission with Direct Temporal Provenance}

Each semantically enriched chunk is written as a Markdown file with YAML frontmatter
containing structured metadata:

\begin{lstlisting}[language=YAML, caption={Emitted chunk document with YAML frontmatter}]
---
source_file: corpus/enriched_full.txt
timestamp: 1666-09-02T08:00
topic: domestic
context: Home
chunk_index: 0
---
[chunk text content]
\end{lstlisting}

\textbf{Critical innovation: direct temporal database writes.} The timestamp in the
frontmatter is written directly from the parsed ISO timestamp in the source format. No
language model is involved in temporal grounding at any stage. This completely bypasses
the temporal extraction failure mode that plagues small-parameter local models, making
the temporal dimension fully reliable regardless of model capability. The temporal
provenance is a first-class metadata field on every chunk, available for filtering,
sorting, and graph construction in all downstream stages.

% -----------------------------------------------------------
\subsection*{4.\ \ Stage 3: Structural Knowledge Graph Construction}

The system delegates to a structural graph builder that indexes all emitted chunk files
into a SQLite database. The structural index provides:

\textbf{BM25 Lexical Search.} Full-text search over chunk content using the Okapi BM25
ranking function. This supports exact term matching, phrase queries, and Boolean
combinations.

\textbf{Typed Graph Edges.} The system constructs typed edges between nodes:

\begin{itemize}[leftmargin=2cm, itemsep=3pt]
  \item CONTAINS edges linking source documents to their constituent chunks;
  \item SIMILAR edges linking chunks with high lexical overlap;
  \item TEMPORAL edges linking chunks that are temporally adjacent in the source corpus.
\end{itemize}

\textbf{Metadata Filtering.} All frontmatter fields (timestamp, topic, context, source
file, chunk index) are indexed as filterable metadata columns, enabling efficient
predicate pushdown for combined structural and metadata queries.

The validated structural graph for the reference corpus contains 29,402 nodes and
355,250 edges in a 109 MB SQLite database.

% -----------------------------------------------------------
\subsection*{5.\ \ Stage 4: Semantic Vector Index Construction}

The system builds a semantic vector index using dense embeddings from a pre-trained
sentence transformer model. In the preferred embodiment, the model is
\texttt{all-mpnet-base-v2}, producing 768-dimensional float32 vectors.

The vector index is stored in a columnar vector database (LanceDB in the preferred
embodiment) optimized for approximate nearest-neighbor (ANN) search. The index supports
cosine similarity queries with sub-millisecond latency.

\textbf{Critical design point:} All embedding computation occurs during this build
stage. At query time, the vector index serves pre-computed embeddings directly. No
model is loaded and no inference is performed to answer a query.

The validated vector index for the reference corpus occupies 102 MB for 6,647 chunk
embeddings at 768 dimensions.

% -----------------------------------------------------------
\subsection*{6.\ \ Stage 5: Multi-Process Parallel Embedding}

The embedding pipeline operates independently of the knowledge graph build, producing
an embedding cache optimized for analysis:

\begin{lstlisting}[caption={Multi-process parallel embedding pipeline}]
Parse -> (optional temporal sample) -> shard across workers -> each worker:
    load sentence-transformer model locally -> encode shard -> return float32 array
-> concatenate shards in original order -> write cache
\end{lstlisting}

\textbf{Key design decisions providing the inventive step:}

\textbf{Independent model instances per worker.} Each worker process loads its own copy
of the sentence-transformer model. There is no shared memory between workers and no
Global Interpreter Lock (GIL) contention. This is in contrast to conventional
approaches that share a single model instance across threads, which in Python-based
systems creates GIL contention that prevents true parallelism.

\textbf{Spawn start method enforcement.} On POSIX systems, the system enforces the
\texttt{spawn} start method for subprocess creation rather than the default \texttt{fork}
method. This ensures clean subprocess isolation, preventing issues with inherited file
descriptors, CUDA contexts, or other process state that can cause subtle failures in
forked multiprocessing.

\textbf{Temporally-uniform sampling.} When processing a subset of the corpus, temporal
sampling uses uniform stride across the time axis rather than random sampling, ensuring
that the embedding cache provides representative manifold coverage of the full temporal
span.

\textbf{Aligned output format.} The output consists of three aligned arrays ---
embeddings, texts, and timestamps --- stored as JSON. This self-contained format
requires no database dependency and is portable across environments.

The architecture achieves near-linear scaling with CPU count on encode-bound workloads.
In the reference implementation, four worker processes embed 6,647 chunks as part of a
full pipeline execution completing in approximately three minutes.

% -----------------------------------------------------------
\subsection*{7.\ \ Stage 6: Embedding Cache}

The JSON embedding cache stores three parallel arrays:

\begin{lstlisting}[language=JSON, caption={Embedding cache JSON structure}]
{
  "embeddings": [[float32, ...], ...],
  "texts": ["chunk text", ...],
  "timestamps": ["ISO timestamp", ...]
}
\end{lstlisting}

The arrays are index-aligned: \texttt{embeddings[i]}, \texttt{texts[i]}, and
\texttt{timestamps[i]} all correspond to the same source chunk. This format is
intentionally simple, requiring no database dependency and supporting direct consumption
by any numerical computing environment.

% -----------------------------------------------------------
\subsection*{8.\ \ Unified Pipeline Orchestration with Snapshot System}

The system provides a unified orchestrator that executes the full pipeline --- NLP
transformation, corpus emission, structural graph indexing, vector index construction,
and snapshot capture --- through a single invocation.

\textbf{Point-in-time snapshots} capture corpus metrics (entry count, chunk count,
temporal span), graph metrics (node count, edge count, index sizes), and build metadata
(timestamp, hardware, duration). Snapshots are versioned and enable comparison across
corpus rebuilds, supporting reproducibility and regression detection.

% -----------------------------------------------------------
\subsection*{9.\ \ Hybrid Retrieval Without Query-Time Inference}

The completed knowledge graph supports hybrid retrieval combining:

\begin{itemize}[leftmargin=2cm, itemsep=3pt]
  \item \textbf{BM25 lexical search} for exact term matching and Boolean queries;
  \item \textbf{Vector ANN search} for semantic similarity queries;
  \item \textbf{Metadata filtering} for temporal range queries, topic filtering, and
        context filtering;
  \item \textbf{Graph traversal} for structural navigation (containment, similarity,
        temporal adjacency).
\end{itemize}

All four retrieval modalities operate on pre-computed indices. No model is loaded at
query time. No inference is performed. Query latency is sub-millisecond for the
reference corpus of 29,402 nodes.

\bigskip
\noindent The following table summarizes the validated reference corpus metrics:

\vspace{6pt}
\begin{center}
\begin{tabular}{lrl}
  \toprule
  \textbf{Metric} & \textbf{Value} & \textbf{Notes} \\
  \midrule
  Source entries          & 6,450     & Diary entries, 1660--1669 \\
  Enriched lines          & 23,235    & Pipe-delimited format \\
  Source file size        & 3.3 MB    & \\
  Chunk count             & 6,647     & Sentence-group chunking \\
  Graph nodes             & 29,402    & \\
  Graph edges             & 355,250   & CONTAINS, SIMILAR, TEMPORAL \\
  SQLite database size    & 109 MB    & Structural index \\
  LanceDB vector index    & 102 MB    & 768-dimensional float32 \\
  Combined KG footprint   & $\sim$241 MB & \\
  Embedding dimensions    & 768       & \texttt{all-mpnet-base-v2} \\
  Full pipeline duration  & $\sim$3 min & Apple Silicon, no GPU \\
  Query latency           & $<$1 ms   & Pre-computed indices \\
  \bottomrule
\end{tabular}
\end{center}

% ============================================================
\section*{Claims}

\subsection*{Independent Claims}

\begin{claims}

% --- Claim 1 ------------------------------------------------
\item A computer-implemented method for constructing a queryable knowledge graph from a
timestamped text corpus, comprising:

\begin{claimelems}
  \item receiving a plurality of text entries, each text entry comprising a timestamp,
        a semantic type label, a context category, and natural language content;

  \item for each text entry, applying a multi-phase natural language processing
        transformation comprising:
    \begin{claimelems}
      \item parsing and validating the timestamp;
      \item segmenting the content into one or more chunks according to a configurable
            chunking strategy;
      \item classifying each chunk with one or more topic labels using a configurable
            topic vocabulary;
      \item classifying each chunk with a single context label using a configurable
            context vocabulary;
    \end{claimelems}

  \item emitting each chunk as a structured document file comprising the chunk text and
        metadata frontmatter, wherein the metadata frontmatter includes a temporal
        provenance field written directly from the parsed timestamp of step (b)(i)
        without invoking any language model for temporal extraction;

  \item constructing a structural graph index from the emitted document files, the
        structural graph index comprising nodes representing chunks, typed edges
        representing containment, similarity, and temporal adjacency relationships,
        and a full-text search index;

  \item constructing a semantic vector index by computing dense vector embeddings for
        each chunk using a pre-trained sentence-transformer model and storing the
        embeddings in a columnar vector store supporting approximate nearest-neighbor
        search;
\end{claimelems}

\noindent wherein steps (d) and (e) collectively constitute ingest-time pre-computation
such that all subsequent queries against the knowledge graph are served by vector lookup
and graph traversal without loading any language model or performing any neural network
inference at query time.

% --- Claim 2 ------------------------------------------------
\item A computer-implemented method for parallel vector embedding of a text corpus
without shared state, comprising:

\begin{claimelems}
  \item receiving a plurality of text chunks derived from a timestamped text corpus;

  \item partitioning the text chunks into a plurality of shards, one shard per
        available CPU core;

  \item spawning a plurality of worker processes using a spawn start method, wherein
        each worker process:
    \begin{claimelems}
      \item loads an independent instance of a sentence-transformer model into its own
            process memory;
      \item encodes its assigned shard of text chunks into dense vector embeddings;
      \item returns the embeddings as a contiguous float32 array;
    \end{claimelems}

  \item concatenating the returned embeddings from all worker processes in the original
        corpus order;
\end{claimelems}

\noindent wherein no shared memory is used between worker processes, no Global
Interpreter Lock contention occurs, and the method achieves near-linear scaling with
the number of available CPU cores.

% --- Claim 3 ------------------------------------------------
\item A system for offline semantic pre-computation of a timestamped text corpus,
comprising:

\begin{claimelems}
  \item a multi-phase NLP transformer module configured to parse timestamped entries,
        segment content into chunks, and classify each chunk with topic and context
        labels from user-configurable YAML vocabularies;

  \item a corpus emitter module configured to write each chunk as a structured document
        file with metadata frontmatter containing temporal provenance written directly
        from source timestamps without language model extraction;

  \item a structural graph builder module configured to index the document files into a
        relational database comprising nodes, typed edges, and a BM25 full-text search
        index;

  \item a semantic vector index builder module configured to compute dense vector
        embeddings using a pre-trained sentence-transformer model and store the
        embeddings in a columnar vector database;

  \item a multi-process parallel embedder module configured to shard the corpus across
        worker processes, each worker loading an independent model instance with no
        shared state;

  \item a unified orchestrator configured to execute modules (a) through (e) in
        sequence through a single invocation and to capture point-in-time snapshots of
        corpus and graph metrics;
\end{claimelems}

\noindent wherein the system is configured to execute entirely on consumer-grade CPU
hardware without GPU acceleration and without network connectivity, and wherein all
queries against the constructed knowledge graph are served without loading any language
model or performing any inference.

\end{claims}

% -----------------------------------------------------------
\subsection*{Dependent Claims}

\begin{claims}
\setcounter{claimnum}{3}

% --- Claim 4 ------------------------------------------------
\item The method of Claim~1, wherein the configurable chunking strategy of step
(b)(ii) comprises sentence-group chunking with a configurable number of sentences per
chunk and a configurable maximum character count per chunk.

% --- Claim 5 ------------------------------------------------
\item The method of Claim~1, wherein the multi-label topic classification of step
(b)(iii) uses TF-IDF weighted keyword matching against a YAML vocabulary file that maps
topic labels to keyword lists, and wherein each chunk may receive zero, one, or multiple
topic labels.

% --- Claim 6 ------------------------------------------------
\item The method of Claim~1, further comprising applying temporally-diverse sampling
when processing a subset of the corpus, wherein the sampling uses uniform stride across
the time axis of the corpus to ensure representative temporal coverage.

% --- Claim 7 ------------------------------------------------
\item The method of Claim~1, wherein the structural graph index of step (d) comprises:
\begin{itemize}[leftmargin=1.8cm, itemsep=3pt]
  \item CONTAINS edges linking source document nodes to constituent chunk nodes;
  \item SIMILAR edges linking chunk nodes having lexical overlap exceeding a configurable
        threshold;
  \item TEMPORAL edges linking chunk nodes that are temporally adjacent in the source
        corpus.
\end{itemize}

% --- Claim 8 ------------------------------------------------
\item The method of Claim~1, wherein the semantic vector index of step (e) uses a
768-dimensional embedding space generated by a pre-trained sentence-transformer model,
and wherein the columnar vector store supports cosine similarity approximate
nearest-neighbor queries with sub-millisecond latency.

% --- Claim 9 ------------------------------------------------
\item The method of Claim~2, wherein the spawn start method is enforced on POSIX
systems to ensure clean subprocess isolation and to prevent inheritance of file
descriptors, GPU contexts, or other parent process state.

% --- Claim 10 -----------------------------------------------
\item The method of Claim~2, further comprising, prior to step (b), applying
temporally-uniform sampling using stride-based selection across the time axis of the
corpus.

% --- Claim 11 -----------------------------------------------
\item The system of Claim~3, further comprising a snapshot module configured to capture
point-in-time metrics including entry count, chunk count, temporal span, node count,
edge count, index sizes, and build metadata, and to support comparison between snapshots
captured at different build times.

% --- Claim 12 -----------------------------------------------
\item The system of Claim~3, wherein the structural graph builder module stores the
relational database as a SQLite file and the semantic vector index builder module stores
the vector database as a LanceDB directory, and wherein the combined knowledge graph
footprint for a corpus of approximately 6,500 entries is approximately 241 megabytes.

% --- Claim 13 -----------------------------------------------
\item The method of Claim~1, wherein the knowledge graph supports hybrid retrieval
combining BM25 lexical search, vector approximate nearest-neighbor search, metadata
filtering on temporal range and topic and context labels, and graph traversal across
typed edges, all without loading any model or performing any inference at query time.

% --- Claim 14 -----------------------------------------------
\item The method of Claim~1, wherein the entire pipeline from raw corpus ingestion to
completed knowledge graph executes in less than five minutes for a corpus of up to
10,000 entries on consumer-grade hardware without GPU acceleration.

% --- Claim 15 -----------------------------------------------
\item The system of Claim~3, wherein the user-configurable YAML vocabularies of module
(a) enable domain adaptation to arbitrary timestamped text corpora --- including
personal journals, clinical notes, legal case files, corporate meeting transcripts,
research literature, and biographical archives --- without modification to any system
module.

\end{claims}

% ============================================================
\section*{Abstract}

A system and method for constructing a semantically enriched, temporally grounded
knowledge graph from timestamped text corpora through offline pre-computation. The
system applies multi-phase NLP transformation --- parsing, chunking, multi-label topic
classification, and context classification using domain-configurable YAML vocabularies
--- to produce structured document chunks with temporal provenance written directly from
source timestamps without any language model extraction. The chunks are indexed into a
hybrid knowledge graph comprising a structural relational database with BM25 full-text
search and typed graph edges, and a semantic vector index with dense embeddings from a
pre-trained sentence-transformer model. A multi-process parallel embedding pipeline
shards the corpus across CPU workers, each loading an independent model instance with
no shared state, achieving near-linear scaling. A unified orchestrator executes the
full pipeline through a single invocation with point-in-time snapshot capture. The
resulting knowledge graph serves all subsequent queries --- lexical, semantic, metadata,
and structural --- through pre-computed indices with sub-millisecond latency, no model
loading, and no inference at query time. The system executes entirely on consumer-grade
CPU hardware without GPU acceleration or network connectivity.

\vspace{24pt}
\noindent\rule{\linewidth}{0.4pt}
\vspace{6pt}

\begin{center}
  \small
  \textit{Patent Application -- Eric G. Suchanek, PhD -- Flux-Frontiers -- March 26, 2026}
\end{center}

\end{document}
